% Figure from MATH 590 notes, section 26. Compile with the notes preamble.
\begin{tikzpicture}[scale=1.25]
  % left: B^2 with the straight-line contraction of the boundary loop toward b_0
  \fill[black!6] (0,0) circle (1);
  \foreach \t/\c in {0.2/75, 0.4/60, 0.6/45, 0.8/30}
    \draw[figB!\c, thin] (\t,0) circle ({1-\t});
  \draw[figB, thick] (0,0) circle (1);
  \draw[figB, thick, -{Stealth[length=5pt]}] (88:1) arc (88:92:1);
  \fill[figB] (1,0) circle (1.4pt) node[right, font=\scriptsize, black] {$b_0$};
  \node[font=\small, figB] at (-0.95,0.95) {$f$};
  \node[font=\scriptsize] at (0,-1.3) {in $B^2$: $f \simeq_p e_{b_0}$};
  % the would-be retraction
  \draw[figR, thick, dashed, -{Stealth[length=6pt]}] (1.55,0.25) to[bend left=18] node[above, font=\scriptsize] {$r$\,?} (2.65,0.25);
  % right: S^1, where f generates pi_1
  \begin{scope}[xshift=4.05cm]
    \draw[figB, thick] (0,0) circle (1);
    \draw[figB, thick, -{Stealth[length=5pt]}] (88:1) arc (88:92:1);
    \fill[figB] (1,0) circle (1.4pt) node[right, font=\scriptsize, black] {$b_0$};
    \node[font=\small, figB] at (-0.95,0.95) {$f$};
    \node[font=\scriptsize, figR] at (0,0.12) {$r \circ (\text{contraction})$};
    \node[font=\scriptsize, figR] at (0,-0.2) {cannot exist};
    \node[font=\scriptsize] at (0,-1.3) {in $S^1$: $[f] = 1 \neq 0$ in $\mathbb{Z}$};
  \end{scope}
\end{tikzpicture}
